
Als junior systeembeheerder krijg je vaak de taak om handmatige handelingen te automatiseren. Condities (\textquote{if}, \textquote{elseif}, \textquote{else}) zijn hierin onmisbaar, omdat je script beslissingen moet kunnen nemen op basis van data of invoer.

Jou opdracht is om een PowerShell-script te schrijven dat de status van een lokale of virtuele server controleert en op basis daarvan actie onderneemt. Het te maken script krijgt de naam \textquote{Check-ServerHealth.ps1}.

De eisen waaraan het moet voldoen zijn:
\begin{enumerate}
\item Vrije schijfruimte controleren
	\begin{itemize}
	\item Vraag de gebruiker via `Read-Host` om het percentage vrije schijfruimte op te geven (een getal tussen 0 en 100).
	\item Schrijf een `if`-statement:
	\item Als het percentage kleiner is dan 10, verschijnt er een kritieke waarschuwing in het rood.
	\item Als het percentage kleiner is dan 25, verschijnt er een gele waarschuwing.
	\item Anders is de schijfruimte in orde en verschijnt er een groene melding.
	\end{itemize}
\item Uitvoeren op disks
	\begin{itemize}
	\item Herschrijf het script en vraag de gebruiker om de naam van een disk (C of D)
	\item Controleer of de variabele niet leeg is (null-check) voordat je ermee gaat werken
	\item Haal de vrije schijfruimte op van de disk (Get-Volume -DriveLetter <driveletter>)
	\item Gebruik dezelfde test(s) als in de vorige opdracht
	\end{itemize}
\item Service-status controleren met logische operators
	\begin{itemize}
	\item Maak twee variabelen aan in je script, bijvoorbeeld:
	\item Vraag aan de gebruiker om een service naam
	\item Vraag aan de gebruiker of er een maintenance window is. De enige geldige antwoorden zijn J of N.
	\item Controlleer of de service draait
	\item Schrijf een if conditie waarin je controleert of een service draait \textbf{en} of er een onderhoudsvenster actief is:
		\begin{itemize}
		\item Als de service niet draait \textbf{en} er is geen onderhoudsvenster, bericht dit dan aan de gebruiker
		\item Gebruik in alle andere gevallen een melding dat er geen actie vereist is.
		\end{itemize}
	\end{itemize}
\end{enumerate}

